\chapter{The Opaque Regime: Deal Marks and Macro Factors}
\label{ch:us}

Part~I assumed that fundamentals arrive on a statutory clock: every firm, every fiscal year,
a filed income statement. In the United States that clock does not exist. What the vendor
delivers instead is the firm's \emph{transaction record}: a sequence of deal events, each
carrying marks --- a valuation, a round classification, an employee count, occasionally a
revenue or EBITDA figure. Between events we learn nothing. This chapter re-derives the
Part~I pipeline under that observation model. Three things change: the \emph{target} (EBITDA
growth must be replaced by proxies constructed from marks), the \emph{unit of analysis} (the
firm-year panel row is replaced by the consecutive-deal pair), and --- most consequentially
--- the \emph{bias structure}: outcomes are observed only when the firm transacts again, so
every naive estimate is filtered through survival-to-next-round. The repair for that
selection is a model of the deal-arrival process itself, which is why Part~III is not an
appendix to this chapter but its completion.

\section{The observation model}
\label{sec:us-obs}

Let $S_{i,t} \in \mathbb{R}^d$ denote the latent state of firm $i$: fundamentals (revenue,
margin trajectory, unit economics) plus soft state (team quality, runway). In the EU regime
we observed a noisy annual census of $S$. Here we observe marks only at the firm's own deal
times $T_{i,1} < T_{i,2} < \cdots$:
\begin{equation}
m_{i,k} \;=\; \big(V_{i,k},\ \texttt{round\_type}_{i,k},\ E_{i,k},\ R_{i,k},\ B_{i,k}\big),
\label{eq:us-marks}
\end{equation}
where $V$ is post-money valuation (populated only for priced rounds), $E$ the reported
employee count, and $R$, $B$ the revenue and EBITDA fields --- populated for roughly $45\%$
and $25\%$ of deal rows respectively in our synthetic mirror, calibrated to vendor coverage
(Chapter~\ref{ch:data}). The marks are noisy, truncated views of $S_{i,T_{i,k}}$: valuation
confounds fundamentals with the market multiple, employees confound productivity with cash
deployment, and the revenue field is populated selectively.

\begin{definition}[LOCF panel]
\label{def:us-locf}
Let $K_i(t) = \max\{k : T_{i,k} \le t\}$, with $K_i(t) = 0$ before the first deal. The
\emph{last-observation-carried-forward panel} is $\tilde{S}_{i,t} = m_{i,K_i(t)}$ together
with the staleness clock $A_{i,t} = t - T_{i,K_i(t)}$. Every firm-side feature in this
chapter is a function of $(\tilde{S}_{i,t}, A_{i,t})$ and the full past mark history, never
of marks dated after $t$.
\end{definition}

The staleness clock is not decoration: a valuation carried forward $4$ years is a different
object from one carried forward $4$ months, and models must be allowed to discount stale
marks. But the deeper problem is \emph{when} the panel refreshes.

\begin{warning}[Sampling times are informative]
\label{warn:us-mnar}
The observation times $T_{i,k}$ are not an exogenous measurement schedule; they are outcomes
of the firm's own trajectory. A firm raises when it is doing well (an up round at a rich
multiple) or when it must (runway exhaustion, a down round); firms that are neither raise
nothing and go dark. Formally, the event indicator $\dd N_i(t)$ is not independent of
$S_{i,t}$ given the past: the missingness mechanism of the LOCF panel is
\emph{missing-not-at-random}. Any panel method that treats $\tilde{S}_{i,t}$ as if generated
by random inspection times --- imputation, mixed models, naive regressions on the LOCF
columns --- inherits a bias whose sign and size depend on the unmodelled sampling mechanism.
The only honest formulations either (i) condition explicitly on the event frame
(Section~\ref{sec:us-event}) and correct for selection (Section~\ref{sec:us-selection}), or
(ii) model the sampling process jointly with the state, which is Part~III.
\end{warning}

\section{What replaces EBITDA growth}
\label{sec:us-proxies}

The EU-track success labels of Chapter~\ref{ch:success} are defined on filed fundamentals.
We now construct their US surrogates, all defined on a pair of consecutive deals
$(T_{i,k}, T_{i,k+1})$ with gap $\Delta_{i,k} = T_{i,k+1} - T_{i,k}$.

\paragraph{(a) Implied valuation growth.}
When both endpoints are priced rounds,
\begin{equation}
g^{\mathrm{val}}_{i,k} \;=\; \frac{\log V_{i,k+1} - \log V_{i,k}}{\Delta_{i,k}},
\label{eq:us-gval}
\end{equation}
an annualized log growth rate. Writing valuation as fundamentals times a sector multiple,
$\log V_{i,k} = \log F_{i,T_{i,k}} + \log m_{s(i),T_{i,k}} + \varepsilon_{i,k}$, gives the
decomposition
\begin{equation}
g^{\mathrm{val}}_{i,k}
\;=\; \underbrace{g^{F}_{i,k}}_{\text{fundamentals growth}}
\;+\; \underbrace{g^{m}_{i,k}}_{\text{multiple expansion}}
\;+\; \underbrace{\nu_{i,k}}_{\text{pricing noise}} .
\label{eq:us-decomp}
\end{equation}
Only $g^{F}$ is a success signal; $g^{m}$ is a market-wide repricing that the firm did not
earn. We identify the multiple component by the public-market-comparables logic: regress
$g^{\mathrm{val}}$ on sector $\times$ time-bucket fixed effects plus macro controls (rate
changes, public-index sector returns over the gap); the fitted cell means estimate
$g^{m}$, and the residual is the firm-level proxy $\hat{g}^{F}_{i,k}$. The fixed effects
must be estimated on all pairs in the cell, not per firm, or they absorb the signal.

\paragraph{(b) Round direction as an ordinal outcome.}
The field \texttt{vc\_round\_up\_down\_flat} classifies each round against the previous one
as $\{\text{down}, \text{flat}, \text{up}\}$. It is coarse but far better populated than
paired valuations (it is defined even when one endpoint is unpriced) and robust to
level noise in $V$. We treat it as an ordinal target (proportional-odds or ordinal
gradient boosting); its drawback is the vendor's flat band, which is a definitional
artifact rather than an economic category.

\paragraph{(c) Employee growth.}
$g^{\mathrm{emp}}_{i,k} = (\log E_{i,k+1} - \log E_{i,k})/\Delta_{i,k}$ is a real-activity
proxy available for most pairs. Its bias is mechanical: headcount growth immediately after
a raise reflects cash deployment, not demonstrated productivity, so it leads fundamentals
optimistically by construction.

\paragraph{(d) Reported revenue growth.}
\texttt{revenue\_growth\_since\_last\_deal} is the most direct surrogate when populated,
but with $\sim\!45\%$ endpoint coverage the pair-level coverage is roughly $20\%$, and
reporters skew to larger, later-stage firms --- the missingness is again not at random.

\begin{table}[htbp]
\centering
\small
\begin{tabular}{p{3.1cm}p{2.1cm}p{2.5cm}p{4.2cm}}
\toprule
Proxy & Pair coverage & Noise & Dominant bias \\
\midrule
$g^{\mathrm{val}}$ \eqref{eq:us-gval} & moderate (both rounds priced) & high (negotiation,
multiple) & upward; confounds multiple expansion unless decomposed via
\eqref{eq:us-decomp} \\
\texttt{vc\_round\_}\newline\texttt{up\_down\_flat} & high & low but coarse & flat band
definitional; censors magnitude \\
$g^{\mathrm{emp}}$ & high & moderate & upward after raises (cash-fueled hiring) \\
\texttt{revenue\_growth\_}\newline\texttt{since\_last\_deal} & low ($\sim\!20\%$) & low
given reported & reporters skew large/late-stage; MNAR \\
\bottomrule
\end{tabular}
\caption{Growth proxies on the consecutive-deal pair frame. All four are additionally
subject to the survival-to-next-round selection of Section~\ref{sec:us-selection}, which
no choice among them avoids.}
\label{tab:us-proxies}
\end{table}

In practice we fit all four and read them jointly: agreement across proxies with different
bias directions is the closest available substitute for a filed income statement.

\section{Supervised learning on the event frame}
\label{sec:us-event}

The prediction unit is the ordered pair $(i,k) \to (i,k+1)$: one row per consecutive-deal
pair, outcome one of the proxies of Section~\ref{sec:us-proxies}. Features come in two
blocks with different epistemic status.

\emph{Firm-history block}, measurable at $T_{i,k}$: the full mark history through the
$k$-th deal (valuation path, round sequence, employee path, deal sizes, investor counts),
the staleness pattern, firm age, sector, and gap lengths between \emph{past} deals.
Nothing dated inside $(T_{i,k}, T_{i,k+1}]$ from the firm's own record enters --- there is
nothing there to use except the second deal itself, which is the outcome's timestamp.

\emph{Gap macro block}, aggregates of exogenous series over the gap: mean policy rate,
change in CPI year-on-year, change in unemployment, and sector deal heat (sector deal
counts per week, normalized by trailing average), all built from point-in-time vintages as
in Chapter~\ref{ch:data}. Because these are external to the single firm, conditioning on
them does not leak the firm's outcome; but it does change the estimand. With the gap block
included, the regression is an \emph{attribution} model --- growth conditional on the
realized macro path and the realized gap length. For deployment as a forecast at $T_{i,k}$,
the gap block must be replaced by scenario paths or integrated over a gap-length model;
the gap length $\Delta_{i,k}$ is itself informative (short gaps select hot firms), and we
return to precisely this point below.

The learner is the Chapter~\ref{ch:xgb} stack unchanged: gradient boosting with blocked
cross-validation (folds split by firm \emph{and} time), early stopping on the temporal
holdout, and conformalized quantile intervals for uncertainty; we do not repeat that
machinery here. Two adaptations matter. First, conformal calibration sets must be blocked
by firm, since pairs from one firm share latent quality and violate exchangeability across
naive random splits. Second, once the selection weights of the next section are
introduced, calibration must be weight-aware (weighted conformal quantiles), otherwise
coverage holds for the selected population only.

\section{The selection problem}
\label{sec:us-selection}

This section is the core of the chapter. A pair $(i,k)$ exists in the training frame only
if the firm raised again within the observation window: define
\begin{equation}
D_{i,k} \;=\; \ind{T_{i,k+1} \le \tau_{\mathrm{end}}},
\label{eq:us-selind}
\end{equation}
where $\tau_{\mathrm{end}}$ is the data cutoff. Every proxy of
Section~\ref{sec:us-proxies} is observed conditional on $D_{i,k} = 1$: we see the growth
of firms \emph{that came back}. The observed distribution is $P(g \mid D{=}1)$, not
$P(g)$, and the gap between them has an exact expression:
\begin{equation}
\E[g \mid D{=}1] - \E[g]
\;=\; \frac{\Cov\!\big(g,\ \ind{D{=}1}\big)}{\Prob(D{=}1)} .
\label{eq:us-selbias}
\end{equation}
Both selection channels have positive covariance with \emph{appearing} in the frame but
opposite implications for the observed growth mix: winners raise again (pulling observed
$g$ up), and cash-constrained firms raise again at down rounds (pulling it down), while
the quiet failures --- firms that never transact again --- contribute nothing at all. The
left tail of $P(g)$ is therefore systematically amputated: shutdown is the one outcome
that generates no second deal. Estimates of mean growth, of feature effects, and of
ranking quality are all computed on the surviving subpopulation.

\paragraph{Bounds.}
If $g$ is bounded to $[g_{\mathrm{lo}}, g_{\mathrm{hi}}]$ (after winsorization, or by
construction for the ordinal outcome), then with $p = \Prob(D{=}1)$,
\begin{equation}
\E[g] \;\in\; \big[\, p\,\E[g \mid D{=}1] + (1-p)\,g_{\mathrm{lo}},\;\;
p\,\E[g \mid D{=}1] + (1-p)\,g_{\mathrm{hi}} \,\big],
\label{eq:us-bounds}
\end{equation}
the worst-case (no-assumption) bounds. With $p$ well below one --- and
Figure~\ref{fig:us-deals-per-firm} shows how few firms generate long deal sequences ---
these are wide. They are still worth reporting: they are the honest statement of what the
data alone say, and any point estimate inside them is purchased with an assumption.

\begin{assumption}[Selection on observed history]
\label{ass:us-soh}
$g_{i,k} \indep D_{i,k} \mid \Hist_{T_{i,k}}$, where $\Hist_{T_{i,k}}$ collects the firm's
mark history through deal $k$ and the macro state: given everything observed up to the
first deal of the pair, whether the firm raises again carries no further information about
its growth.
\end{assumption}

\paragraph{The practical fix: inverse-probability-of-observation weights.}
Under Assumption~\ref{ass:us-soh}, define the observation probability
\begin{equation}
\pi_{i,k} \;=\; \Prob\big(D_{i,k}{=}1 \mid \Hist_{T_{i,k}}\big)
\;=\; 1 - \exp\!\Big(-\!\int_{T_{i,k}}^{\tau_{\mathrm{end}}} \lambda_i(u)\dd u\Big),
\label{eq:us-pi}
\end{equation}
where $\lambda_i$ is the firm's next-deal intensity given its history. Weighting each
observed pair by $1/\hat{\pi}_{i,k}$ (stabilized and trimmed in the usual way) makes the
weighted frame representative of all at-risk firms: $\E[D g / \pi] = \E[g]$. This is
inverse-probability weighting for informative censoring in the sense of
\cite{hernan2020}, with the hazard-model lineage of \cite{cox1972}; validity requires
$\pi_{i,k} > 0$ and a well-specified intensity, and residual selection on latent quality
beyond $\Hist$ is not repaired --- which is why we report \eqref{eq:us-bounds} alongside.

\begin{remark}[The bridge to Part~III]
\label{rem:us-bridge}
Equation~\eqref{eq:us-pi} is the structural link of this manual. The quantity
$\lambda_i(u)$ --- the conditional intensity of the firm's deal-arrival process --- is
exactly what Part~III calibrates: inhomogeneous Poisson with macro covariates
(Chapter~\ref{ch:poisson}), self-exciting and inhibiting extensions
(Chapter~\ref{ch:hawkes}), and censored variants (Chapter~\ref{ch:censored}). The count
model is therefore not a separate deliverable that happens to share the data: it
\emph{repairs the growth regression}. A team that skips the point-process layer does not
merely lose the deal-timing forecast; it loses the only principled weights for its growth
model. Chapter~\ref{ch:selection} validates the combined pipeline --- intensity fit,
weights, weighted growth regression --- on synthetic ground truth where $P(g)$ is known.
\end{remark}

\begin{figure}[htbp]
\centering
\includegraphics[width=.85\textwidth]{fig_deals_per_firm}
\caption{Rank-size plot of deals per firm in both synthetic regimes. The event record is
heavily concentrated: a small minority of serially raising firms contributes most
consecutive-deal pairs, while the modal firm contributes one deal and zero pairs. The
pair frame therefore oversamples exactly the firms with the highest historical deal
hazard; unweighted estimates describe serial raisers, not the tracked universe. Firm-level
cluster bootstrap is mandatory for all standard errors on this frame.}
\label{fig:us-deals-per-firm}
\end{figure}

\section{Macro factors as drivers}
\label{sec:us-macro}

With firm-level outcomes this sparse and selected, much of the macro question --- do rate
changes move private-market activity and pricing? --- is better asked at the sector level.
The workhorse is a sector-exposure (shift-share style) design. Let $y_{s,w}$ be a
sector-week outcome (log deal count, or the cell-mean valuation growth from
\eqref{eq:us-decomp}), $\Delta r_w$ the rate change, and $e_s$ a \emph{predetermined}
exposure measure (share of pre-revenue firms in the sector, a duration proxy for how far
cash flows sit in the future). We fit
\begin{equation}
y_{s,w} \;=\; \alpha_s + \gamma_w + \delta\,(e_s \cdot \Delta r_w) + u_{s,w},
\label{eq:us-shiftshare}
\end{equation}
with sector and time fixed effects and two-way clustered standard errors. The time effect
$\gamma_w$ absorbs every common shock --- including the \emph{average} effect of the rate
change itself --- so $\delta$ is identified purely from differential sensitivity: do
high-exposure sectors move more when rates move?

The interpretation limits are sharp and worth stating. First, the common component of the
macro effect is not identified by design; \eqref{eq:us-shiftshare} estimates relative
sensitivities only. Second, latent risk appetite co-moves with rates; if its sector
loadings correlate with $e_s$, the differential estimate is contaminated too. Third,
exposure must be predetermined and not itself an outcome of recent sector performance. A
design-based reading in the spirit of \cite{imbens2015} is available only under an
explicit exogeneity assumption on the exposure shares, which we state and stress-test on
the testbed rather than assume silently. Where valuation marks are the outcome, all of
Section~\ref{sec:us-selection}'s caveats apply on top.

This is the forward pointer of the chapter: deal \emph{counts} are complete (a deal either
appears in the feed or does not), while marks are sparse and selection-laden. The cleanest
macro inference therefore lives at the count layer --- intensities modulated by macro
covariates --- which is where Chapter~\ref{ch:poisson} begins.

\section{Decision table: EU track versus US track}
\label{sec:us-decision}

\begin{table}[htbp]
\centering
\small
\begin{tabular}{p{2.5cm}p{4.6cm}p{5.6cm}}
\toprule
Aspect & EU track (Part~I) & US adaptation (this chapter) \\
\midrule
Target & EBITDA/revenue growth from filings & proxy set: $\hat{g}^{F}$ from
\eqref{eq:us-decomp}, ordinal round direction, employee growth, reported revenue growth \\
Unit & firm $\times$ fiscal year & consecutive-deal pair $(i,k)$ \\
Feature clock & fiscal-year as-of with publication lags & event-time as-of with staleness
clock $A_{i,t}$ \\
Dominant bias & leakage via publication timing & survival-to-next-round selection; MNAR
sampling times \\
Primary fix & point-in-time joins, blocked CV & IP-of-observation weights from the
next-deal intensity \eqref{eq:us-pi}; worst-case bounds \eqref{eq:us-bounds} \\
Uncertainty & conformal intervals, firm-blocked & weight-aware conformal, firm-cluster
bootstrap \\
Macro layer & covariates in the growth model & shift-share differentials
\eqref{eq:us-shiftshare}; count-layer inference in Part~III \\
\bottomrule
\end{tabular}
\caption{Method mapping between observation regimes. Every row of the US column degrades
gracefully to the EU column when fundamentals become observable; the converse direction is
the content of this chapter.}
\label{tab:us-decision}
\end{table}

The table is the chapter in miniature. Nothing in the US column is exotic; each entry is
the Part~I method re-derived under a harsher observation model, with one genuinely new
ingredient --- the deal-arrival intensity --- that the next three chapters construct.
